% Derived from the local CircuitikZ manual's blockdef style.
\documentclass[border=2pt]{standalone}
\usepackage{circuitikz}

\begin{document}
\tikzset{
  stage/.style={draw,minimum width=2.2cm,minimum height=8mm},
  blockline/.style={blue,line width=0.8pt}
}
\begin{tikzpicture}
  \node[stage] (input) at (0,0) {input stage};
  \node[stage] (output) at (5.2,0) {output stage};
  \draw[-{Straight Barb}] (input.east) -- (output.west);
  \coordinate (block-left) at ([yshift=-8mm]input.south west);
  \coordinate (block-right) at ([yshift=-8mm]output.south east);
  \draw[blockline] (block-left) -- (block-right);
  \draw[blockline,-{Straight Barb[harpoon,reversed,left,length=0.2cm]}]
    (block-left) -- (block-right);
  \draw[blockline,-{Straight Barb[harpoon,reversed,right,length=0.2cm]}]
    (block-right) -- (block-left);
  \path (block-left) -- node[midway,fill=white,inner sep=2pt] {amplifier block} (block-right);
\end{tikzpicture}
\end{document}
